\chapter{Build one Moran event from its biological rules}
\label{chap:moran-event}

\section{Start with a known population state}

We now turn the biological description into explicit probabilities. The
population size is \(N\). The current A count is \(i\), so the current B
count is \(N-i\). Each A individual has reproductive weight \(r=1+s\),
and each B individual has weight \(1\).

One event has two selections. First we choose a parent. Then we independently
choose an individual for replacement from the \emph{current} population.
The parent's offspring has its parent's type and replaces the selected
individual. We use the old counts for both selections.

We are not choosing two distinct individuals without replacement. The parent
can also be selected for replacement. In that case its offspring takes its
place and the type count does not change. Allowing this possibility is part
of the particular Moran model used here.

\section{Choose the parent using reproductive weights}

\subsection{Add the weights within each type}

There are \(i\) A individuals, each with weight \(r\). Their total weight
is \(ri\), which means \(r\) multiplied by \(i\). There are \(N-i\)
B individuals, each with weight one. Their total weight is \(N-i\).

Let \(W_i\) denote the total reproductive weight of the population when
the A count is \(i\). The subscript \(i\) reminds us that the total depends
on the current state. Adding the two type weights gives
\[
  W_i=ri+N-i.
\]

\subsection{Divide a type's weight by the total}

Choosing a parent proportional to weight means that a type's chance is its
total weight divided by the population's total weight:
\begin{align*}
 \Pr(\text{parent is A}\mid i)&=\frac{ri}{W_i},\\
 \Pr(\text{parent is B}\mid i)&=\frac{N-i}{W_i}.
\end{align*}
The notation \(\mid i\) means that these chances are calculated given the
current count \(i\); the fixed values of \(N\) and \(r\) are understood.
The numerators add to \(W_i\), so the probabilities add to one.

\begin{workedexample}
Take \(N=5\), \(i=2\), and \(s=0.5\). Then \(r=1.5\). The two A
individuals have total weight \(1.5\times2=3\). The three B individuals
also have total weight \(1\times3=3\). The total is \(W_i=3+3=6\).

The chance of choosing an A parent is \(3/6=0.5\). The chance of choosing
a B parent is also \(3/6=0.5\). Although there are fewer A individuals,
their greater per-individual weight makes the two types equally likely to
supply the next parent in this particular state.
\end{workedexample}

\section{Choose the individual for replacement uniformly}

The replacement selection does not use reproductive weights. Every current
individual has probability \(1/N\) of being selected. Because there are
\(i\) A individuals and \(N-i\) B individuals,
\begin{align*}
 \Pr(\text{replacement is A}\mid i)&=\frac{i}{N},\\
 \Pr(\text{replacement is B}\mid i)&=\frac{N-i}{N}.
\end{align*}

The phrase ``replacement is A'' describes the type of the old individual
being replaced, not the type of the new offspring.

For the five-individual example, the replacement chances are \(2/5=0.4\)
for A and \(3/5=0.6\) for B. They remain these values regardless of which
type was chosen as the parent. We have not removed anyone between the draws.

\section{List all four parent--replacement combinations}

There are two parent types and two replacement types, giving four possible
combinations. Conditional on the current state, the two draws are independent.
We therefore multiply the appropriate parent and replacement probabilities.

\begin{center}
\begin{tabular}{llll}
\hline
Parent & Old individual replaced & Probability & A-count change \\
\hline
A & A & \(\dfrac{ri}{W_i}\dfrac{i}{N}\) & \(0\) \\[6pt]
A & B & \(\dfrac{ri}{W_i}\dfrac{N-i}{N}\) & \(+1\) \\[6pt]
B & A & \(\dfrac{N-i}{W_i}\dfrac{i}{N}\) & \(-1\) \\[6pt]
B & B & \(\dfrac{N-i}{W_i}\dfrac{N-i}{N}\) & \(0\) \\[6pt]
\hline
\end{tabular}
\end{center}

If A produces an offspring that replaces B, the population gains one A.
If B produces an offspring that replaces A, the population loses one A.
If parent and replacement have the same type, the count is unchanged.

The two unchanged cases include more than replacement of the parent itself.
An A parent can replace a different A individual too. We only record type
counts, so both possibilities have the same effect on our state.

\section{Name the three transition probabilities}

A \textbf{transition} is a move from the current state to the next state.
Let \(T_i^+\) mean the probability of moving from \(i\) to \(i+1\).
Let \(T_i^-\) mean the probability of moving from \(i\) to \(i-1\).
Let \(T_i^0\) mean the probability of remaining at \(i\).

The superscripts \(+\), \(-\), and \(0\) are labels for increase,
decrease, and no change. They are not exponents in this notation.

Reading the four-case table gives
\begin{align}
 T_i^+&=\frac{ri}{ri+N-i}\frac{N-i}{N},\label{eq:t-plus}\\
 T_i^-&=\frac{N-i}{ri+N-i}\frac{i}{N},\label{eq:t-minus}\\
 T_i^0&=1-T_i^+-T_i^-.
 \label{eq:t-stay}
\end{align}

The last equation follows because these three outcomes are mutually exclusive
and cover every possibility. It is also the sum of the A--A and B--B
probabilities in the table.

\subsection{Check that the probabilities make sense}

Every factor in the first two expressions is a nonnegative probability.
The four-case table contains all possible results of the two draws, so its
entries add to one. Therefore \(T_i^0\) is nonnegative too, and
\[
  T_i^++T_i^-+T_i^0=1.
\]

At \(i=0\), the numerator of each change probability contains a factor
\(i=0\). Thus \(T_0^+=T_0^-=0\) and \(T_0^0=1\). At \(i=N\), the
factor \(N-i=0\) gives the same conclusion. The two boundaries really are
absorbing under these rules.

\section{Work through the advantageous example completely}

Return to \(N=5\), \(i=2\), and \(s=0.5\), so \(r=1.5\) and \(W_i=6\).
The parent chances are \(0.5\) for each type. The replacement chances are
\(0.4\) for A and \(0.6\) for B. The four joint probabilities are
\begin{align*}
 \Pr(\text{A parent, A replaced}\mid i=2)&=0.5(0.4)=0.2,\\
 \Pr(\text{A parent, B replaced}\mid i=2)&=0.5(0.6)=0.3,\\
 \Pr(\text{B parent, A replaced}\mid i=2)&=0.5(0.4)=0.2,\\
 \Pr(\text{B parent, B replaced}\mid i=2)&=0.5(0.6)=0.3.
\end{align*}
Consequently,
\[
  T_2^+=0.3,\qquad T_2^-=0.2,\qquad T_2^0=0.2+0.3=0.5.
\]

The expected count change is \(1(0.3)+(-1)(0.2)+0(0.5)=0.1\). Selection
pushes the average upward. The next event can still decrease the count,
and in this state it does so with probability \(0.2\).

\subsection{Follow two particular random events}

One way to perform the two selections is to draw two independent uniform
random numbers. A uniform random number between zero and one assigns a
probability equal to its length to each interval. For example, the interval
from zero to \(0.5\) has probability \(0.5\).

Let \(a\) denote the number used for the parent draw, and let \(b\) denote
the number used for the replacement draw. In this state, choose an A parent
if \(a<0.5\), and choose an A replacement if \(b<0.4\). Otherwise choose B
for the corresponding selection.

\begin{workedexample}
\textbf{First event, starting with two A individuals.} Suppose \(a=0.2\)
and \(b=0.8\). Since \(0.2<0.5\), the parent is A. Since \(0.8\) is not
less than \(0.4\), the old individual selected for replacement is B.
The A offspring replaces B. The count changes from two to three.

\textbf{Second example, resetting to two A individuals.} Suppose \(a=0.8\)
and \(b=0.2\). The parent is B because \(0.8\) is not less than \(0.5\).
The replacement is A because \(0.2<0.4\). The count changes from two to one.

We reset the starting count between these examples. If we continued directly
from the first event instead, we would need to recompute the probabilities
using a count of three.
\end{workedexample}

\subsection{A shorter but equivalent implementation}

Since we only record the count, the program does not need to generate parent
and replacement labels separately. Let \(z\) denote one uniform random
number between zero and one. We can divide its possible range into three
intervals:
\begin{itemize}
 \item If \(z<T_i^+\), increase the count by one.
 \item Otherwise, if \(z<T_i^++T_i^-\), decrease the count by one.
 \item Otherwise, leave the count unchanged.
\end{itemize}
For the example, the interval lengths are \(0.3\), \(0.2\), and \(0.5\).
This reproduces the same three outcome probabilities. It is not changing
the biology of the model; it is sampling the already calculated outcomes
more directly.

\section{Repeat the calculation when selection is absent}

Set \(s=0\), so \(r=1\). The total reproductive weight becomes
\(W_i=i+N-i=N\). Both change probabilities simplify to
\[
  T_i^+=T_i^-=\frac{i(N-i)}{N^2}.
\]
The symbol \(N^2\) means \(N\) multiplied by itself. For \(N=5\) and
\(i=2\), this gives
\[
  T_2^+=T_2^-=\frac{2(5-2)}{5^2}=\frac{6}{25}=0.24,
\]
and
\[
  T_2^0=1-0.24-0.24=0.52.
\]
The expected count change is zero because increases and decreases have equal
probability. Nevertheless, nearly half of all event attempts change the
count in this state. Neutrality means no systematic reproductive advantage;
it does not mean no evolution by drift.

\section{Why we keep unchanged events}

Imagine removing every no-change event from a recorded trajectory. The
sequence of states at the remaining changes could still tell us which
boundary was reached. But the number of recorded steps would no longer be
the number of reproduction and replacement attempts.

The simulation in this project keeps all attempts. A plotted horizontal
segment is a real feature of the recorded count process. It is not an error
or an instruction that failed to run.

\begin{checkpoint}
\textbf{Question.} A neutral population has \(N=4\) and \(i=1\). Work
out both parent probabilities, both replacement probabilities, and the three
transition probabilities. What is the expected count change?

\textbf{Worked answer.} Neutrality means \(r=1\), so the total parent
weight is four. An A parent has probability \(1/4\), and a B parent has
probability \(3/4\). The replacement probabilities are also \(1/4\) for A
and \(3/4\) for B. Therefore
\[
 T_1^+=\frac14\frac34=\frac{3}{16},\qquad
 T_1^-=\frac34\frac14=\frac{3}{16}.
\]
The no-change probability is
\(1-3/16-3/16=10/16\). The expected count change is
\(3/16-3/16=0\). A change is possible despite the zero expectation.
\end{checkpoint}

The next task is to express these calculations as instructions that Julia
can execute. The mathematical inputs are just \(N\), \(i\), and \(s\).
The immediate output is the probability of each possible next count.
